\chapter*{Abstract}

Real-time multi-object tracking systems commonly rely on fixed detector operating parameters and tracker confidence thresholds, which may become unsuitable when scene conditions change or when detectors with different confidence-score distributions are introduced. This thesis addresses these limitations by proposing two complementary, training-free control approaches that improve tracking performance without retraining the detector or modifying the underlying tracking algorithms.

The first approach, Adaptive Control for Multi-Object Tracking (AC-MOT), dynamically adjusts detector confidence thresholds, non-maximum suppression thresholds, and input resolution according to scene complexity and tracking information. The proposed controller employs three operating states with stability and recovery mechanisms, while its parameters are selected using cross-validation and fixed before final evaluation. Experiments on seven VisDrone2019-MOT sequences demonstrate improvements of 2.937 and 1.325 HOTA points and 17.995 and 6.318 MOTA points with ByteTrack and OATrack, respectively. Additional experiments on UAVDT demonstrate improvements in MOTA of 17.601 and 3.267 points without retuning.

The second approach, the Self-Calibrating Control Layer (SCL), addresses incompatibilities between detector confidence-score distributions and tracker thresholds. It introduces a selective, host-aware adaptation mechanism that analyzes recent detection scores, preserves the original tracker configuration when scores are compatible, and otherwise adjusts detection candidates, scores, and thresholds. Experiments involving seven trackers, four detection sources, and three datasets demonstrate improvements under detector--tracker mismatch conditions. In a controlled stress test involving halved detection scores, SCL restores two trackers from zero to approximately 66 HOTA points. Comparisons with adaptive-threshold baselines and ACT further demonstrate improvements in the evaluated configurations.

The proposed controllers introduce low additional computational overhead, averaging approximately 4.2 ms/frame for AC-MOT and 2.95 ms/frame for SCL in their respective experimental settings. The findings demonstrate that adaptive detector operating-point control and detector--tracker interface adaptation can improve tracking performance and robustness without additional model training. The two approaches were evaluated independently, while their joint integration and broader deployment on resource-constrained platforms remain areas for future investigation.

\vspace{0.5cm}
\noindent\textbf{Keywords:} Real-time object detection, multi-object tracking, adaptive control, detector operating point, confidence-score adaptation, UAV tracking, training-free methods.
